\documentclass[12pt]{article}

\usepackage{fontspec}
\newfontfamily\handwriting{a-casual-handwritten-pen}[Path=../assets/fonts/,Extension=.ttf,Scale=1.5]
\newfontfamily\comic{Comic Neue}

\usepackage[utf8]{inputenc}
\usepackage{amsmath}
\usepackage{amssymb}
\usepackage{enumitem}
\usepackage{xcolor}
\usepackage{soul}
\usepackage{tikz}
\usepackage[most]{tcolorbox}
\usepackage{titlesec}
\usepackage{multicol}
\usepackage{environ}

\usetikzlibrary{fit, positioning, arrows.meta}

% --- Custom Colors ---
\definecolor{primaryblue}{RGB}{42, 111, 151}
\definecolor{exbg}{RGB}{255, 240, 245}
\definecolor{rulebg}{RGB}{232, 245, 233}
\definecolor{highlight}{RGB}{255, 243, 176}
\definecolor{pinkanno}{RGB}{255, 105, 180}
\definecolor{purpleanno}{RGB}{147, 112, 219}
\definecolor{solutionbg}{RGB}{245, 245, 255}

% --- Header Styling ---
\titleformat{\section}{\color{primaryblue}\normalfont\Large\bfseries}{\thesection}{1em}{}
\titleformat{\subsection}{\color{primaryblue}\normalfont\large\bfseries}{\thesubsection}{1em}{}

% --- Friendly Boxes ---
\newtcolorbox{friendlyex}[1]{
    colback=exbg,
    colframe=pinkanno!50,
    arc=5pt,
    title=\textbf{#1},
    coltitle=black,
    attach title to upper,
    after title={:\quad}
}

\newtcolorbox{friendlyrule}[1]{
    colback=rulebg,
    colframe=primaryblue!30,
    arc=5pt,
    fonttitle=\bfseries,
    title={#1},
    coltitle=primaryblue
}

\newcommand{\funhl}[1]{\sethlcolor{highlight}\hl{#1}}

% --- Show/Hide Solutions ---
\newif\ifshowsolutions

% Uncomment the next line to show solutions.
\showsolutionstrue

\newcommand{\solutionblank}{\vspace{25ex}}

\NewEnviron{solutionbox}{
\ifshowsolutions
\begin{tcolorbox}[
    colback=solutionbg,
    colframe=purpleanno!45,
    arc=5pt,
    title={Solution},
    coltitle=primaryblue,
    fonttitle=\bfseries
]
\BODY
\end{tcolorbox}
\else
\solutionblank
\fi
}

\begin{document}
\handwriting

\begin{center}
    \Large\textbf{\color{primaryblue}Module 1 Lesson 3\\Linear Functions} \\
    \vspace{0.2cm}
    \hrule
\end{center}

\section*{Objectives}

\begin{friendlyrule}{In this lesson, we will learn how to}
\begin{itemize}
    \item compute the slope of a line through two points;
    \item write and graph a line in slope-intercept form;
    \item compute the average rate of change of a function;
    \item recognize and apply linear functions to model real situations.
\end{itemize}
\end{friendlyrule}

\section*{1. Slope of a Line}

\begin{friendlyrule}{Slope}
\[
m=\frac{\text{rise}}{\text{run}}=\frac{\text{the change in }y}{\text{the change in }x}=\frac{y_2-y_1}{x_2-x_1}.
\]
\end{friendlyrule}

\begin{friendlyex}{Example 1}
Find the slope of the line through each pair of points.

\begin{itemize}
    \item[(1)] $(4,3)$ and $(-2,15)$
    \item[(2)] $(6,-3)$ and $(-2,-3)$
    \item[(3)] $(3,-2)$ and $(3,6)$
\end{itemize}
\end{friendlyex}

 

\begin{solutionbox}
\textbf{(1)}

\[
m=\frac{15-3}{-2-4}=\frac{12}{-6}=-2.
\]

\textbf{(2)}

\[
m=\frac{-3-(-3)}{-2-6}=\frac{0}{-8}=0.
\]

This is a horizontal line.

\textbf{(3)}

\[
m=\frac{6-(-2)}{3-3}=\frac{8}{0},
\]

which is \textbf{undefined}. This is a vertical line.
\end{solutionbox}

  

\section*{2. Slope-Intercept Form}

\begin{friendlyrule}{Slope-Intercept Form}
A line can be represented by an equation of the form

\[
y=mx+b,
\]

where $m$ is the slope and $b$ is the $y$-intercept of the line.
\end{friendlyrule}

For example, the line $y=\dfrac23 x+4$ has slope $\dfrac23$ and $y$-intercept $(0,4)$.

\begin{friendlyex}{Example 2}
\begin{itemize}
    \item[(1)] Given the equation $3x+2y=6$, write the equation in slope-intercept form, and determine the slope and $y$-intercept.
    \item[(2)] Write an equation of a line in slope-intercept form given a point $(2,-4)$ and slope $m=-3$.
\end{itemize}
\end{friendlyex}

 

\begin{solutionbox}
\textbf{(1)} Solve for $y$:

\[
3x+2y=6.
\]

\[
2y=-3x+6.
\]

\[
y=-\frac32 x+3.
\]

The slope is $-\dfrac32$ and the $y$-intercept is $(0,3)$.

\textbf{(2)} Use point-slope form:

\[
y-(-4)=-3(x-2).
\]

\[
y+4=-3x+6.
\]

\[
y=-3x+2.
\]
\end{solutionbox}

  

\section*{3. Average Rate of Change}

\begin{friendlyrule}{Average Rate of Change (AROC)}
The average rate of change of a function over an interval measures the average speed at which the output changes relative to the input.

\[
\text{AROC on }[x_1,x_2]=\frac{f(x_2)-f(x_1)}{x_2-x_1}.
\]
\end{friendlyrule}

\begin{friendlyex}{Example 3}
Given $f(x)=x^2+2$, find the average rate of change on $[1,3]$ and on $[-1,1]$.
\end{friendlyex}

 

\begin{solutionbox}
On $[1,3]$:

\[
\text{AROC}=\frac{f(3)-f(1)}{3-1}=\frac{(9+2)-(1+2)}{2}=\frac{11-3}{2}=4.
\]

On $[-1,1]$:

\[
\text{AROC}=\frac{f(1)-f(-1)}{1-(-1)}=\frac{(1+2)-(1+2)}{2}=\frac{0}{2}=0.
\]
\end{solutionbox}

  

\section*{4. Linear Functions}

\begin{friendlyrule}{Linear Function}
A \funhl{linear function} is a function with a \textbf{constant} average rate of change. Its graph is a straight line, written in slope-intercept form as

\[
f(x)=mx+b,
\]

where $m$ is the rate of change and $b$ is the value of the function when $x=0$.
\end{friendlyrule}

\begin{friendlyex}{Example 4}
Determine which of the following relations represents a linear function.

\[
\begin{array}{c|cccc}
\text{(a)}\ x & 0 & 1 & 2 & 3 \\ \hline
F(x) & 2.5 & 3 & 3.5 & 4
\end{array}
\qquad
\begin{array}{c|cccc}
\text{(b)}\ s & 0 & 2 & 4 & 6 \\ \hline
g(s) & 100 & 90 & 70 & 40
\end{array}
\qquad
\begin{array}{c|cccc}
\text{(c)}\ t & 10 & 20 & 30 & 40 \\ \hline
h(t) & 5.5 & 5 & 4.5 & 4
\end{array}
\]
\end{friendlyex}

 

\begin{solutionbox}
\textbf{(a)} Each step in $x$ increases $F(x)$ by exactly $0.5$. The rate of change is constant, so this \textbf{is} linear.

\textbf{(b)} The changes in $g$ are $-10,-20,-30$ for equal steps of $2$ in $s$, giving rates $-5,-10,-15$. The rate is \textbf{not} constant, so this is \textbf{not} linear.

\textbf{(c)} Each step of $10$ in $t$ decreases $h(t)$ by exactly $0.5$. The rate of change is constant, so this \textbf{is} linear.
\end{solutionbox}

  

\section*{5. Applications of Linear Functions}

\begin{friendlyex}{Example 5}
A city's population was $30{,}700$ in the year 2000 and is growing by $850$ people a year.

\begin{itemize}
    \item[(a)] Give a formula for the population $P$ as a function of the number of years $t$ since 2000.
    \item[(b)] What is the population predicted to be in 2012?
    \item[(c)] When is the population expected to reach $45{,}000$?
\end{itemize}
\end{friendlyex}

 

\begin{solutionbox}
\textbf{(a)}

\[
P(t)=30{,}700+850t.
\]

\textbf{(b)} $2012$ corresponds to $t=12$:

\[
P(12)=30{,}700+850(12)=30{,}700+10{,}200=40{,}900.
\]

\textbf{(c)} Set $P(t)=45{,}000$:

\[
45{,}000=30{,}700+850t.
\]

\[
14{,}300=850t.
\]

\[
t\approx 16.8.
\]

The population is expected to reach $45{,}000$ about $16.8$ years after 2000, i.e., around the year 2017.
\end{solutionbox}

  

\begin{friendlyex}{Example 6}
Annual revenue $R$ (in billions of dollars) from a restaurant chain worldwide can be estimated by

\[
R=19.1+1.8t,
\]

where $t$ is the number of years since January 1, 2005.

\begin{itemize}
    \item[(a)] What is the slope of this function? Interpret it.
    \item[(b)] What is the vertical intercept? Interpret it.
    \item[(c)] What annual revenue does the function predict for 2012?
    \item[(d)] When is annual revenue predicted to hit $35$ billion dollars?
\end{itemize}
\end{friendlyex}

 

\begin{solutionbox}
\textbf{(a)} The slope is $1.8$ billion dollars per year: revenue increases by about $\$1.8$ billion each year.

\textbf{(b)} The vertical intercept is $19.1$ billion dollars (at $t=0$, i.e., 2005): revenue was about $\$19.1$ billion in 2005.

\textbf{(c)} $2012$ corresponds to $t=7$:

\[
R(7)=19.1+1.8(7)=19.1+12.6=31.7\text{ billion dollars.}
\]

\textbf{(d)} Set $R=35$:

\[
35=19.1+1.8t.
\]

\[
15.9=1.8t.
\]

\[
t\approx 8.8.
\]

Revenue is predicted to reach $\$35$ billion about $8.8$ years after 2005, i.e., around late 2013.
\end{solutionbox}

  

\begin{friendlyex}{Example 7}
World grain production was $1241$ million tons in 1975 and $2048$ million tons in 2005, increasing at approximately a constant rate.

\begin{itemize}
    \item[(a)] Find a linear function $P(t)$ for grain production, in million tons, $t$ years since 1975.
    \item[(b)] Interpret the slope.
    \item[(c)] Interpret the vertical intercept.
\end{itemize}
\end{friendlyex}

 

\begin{solutionbox}
\textbf{(a)} The slope is

\[
m=\frac{2048-1241}{2005-1975}=\frac{807}{30}=26.9.
\]

Since $P(0)=1241$,

\[
P(t)=1241+26.9t.
\]

\textbf{(b)} The slope means grain production increased by about $26.9$ million tons each year.

\textbf{(c)} The vertical intercept means grain production was $1241$ million tons in 1975.
\end{solutionbox}

  

\section*{Summary}

\begin{friendlyrule}{Key Ideas}
\begin{itemize}
    \item Slope: $m=\dfrac{y_2-y_1}{x_2-x_1}$. A horizontal line has slope $0$; a vertical line has undefined slope.
    \item Slope-intercept form: $y=mx+b$.
    \item Average rate of change on $[x_1,x_2]$ is $\dfrac{f(x_2)-f(x_1)}{x_2-x_1}$.
    \item A linear function has a constant average rate of change, and its graph is a straight line.
    \item Linear functions are used to model real-world quantities that change at a steady rate.
\end{itemize}
\end{friendlyrule}

\end{document}
